\section{Evaluation Framework}\label{sec:evaluation}

\subsection{Evaluation Design}

Because cluster quality is defined by the oracle, the system cannot be graded against a fixed
target. The evaluation therefore closes a three-party loop: the \textbf{system under test} is
driven by an \textbf{LLM-simulated oracle} that plays a user with a hidden goal, and a separate
\textbf{LLM judge} scores the completed session. The simulated oracle runs through exactly the
same code path as a human client, so its sessions are indistinguishable from live ones in the
database. To avoid a model rewarding outputs that match its own conventions, the system, the
oracle, and the judge each run on a \emph{different} model family.

\subsection{Oracle Formulation}

\paragraph{Ground truth.} Each session is anchored by a \emph{ground truth}: a target navigation
trajectory, an ordered list of operations, paired with a natural-language \emph{intent
description} that paraphrases it. The oracle is shown only the intent description, never the
operation list, so it must reach its goal by stating intent and reacting to what the system
returns, just as a real user would.

\paragraph{Persona.} A persona is layered on top of the ground truth through two dials:
\emph{verbosity}, which controls how much the oracle says, and \emph{patience}, which controls how
many misunderstandings it tolerates before giving up. A persona bundle, the persona dials plus the
ground truth, is a write-once, versioned artefact, so the same experimental condition can be
re-run unchanged.

\paragraph{Protocol.} The oracle follows a strict protocol: one request per message, answer any
clarifying question, ask once per session why a specific film sits in its cluster, and stop either
on success (rating the session 4--5) or on giving up after repeated failures (rating 1--2). A turn
budget caps the session if neither happens.

\paragraph{Reproducibility.} The oracle is reproducible: its output is a deterministic function of
the persona, the ground truth, the session seed, and the turn number, so reruns of the same
quadruple reproduce the same session.

\subsection{Metrics}

Evaluation runs automatically at the end of each session and produces two kinds of measurement.

\paragraph{Deterministic metrics.} These are read directly from the recorded session state and
need no model: the number of clusters in the final snapshot, the number of oracle turns, the
number of navigation operations executed, the fraction of turns on which the clarifier fired, and
the total LLM cost. The oracle's own self-rating (1--5) is recorded alongside them.

\paragraph{Judge dimensions.} These are subjective scores (1--5) assigned by the judge after
reading the full transcript, the per-turn action log, and the final clusters. Four are scored on
every session: \emph{operation appropriateness} (did the system choose the right operation and
parameters for each request), \emph{intent alignment} (does the final result reflect the stated
goal), \emph{label accuracy}, and \emph{clustering coherence}. Three more are conditional, scored
only when the session exercised the relevant feature: \emph{explanation quality}, \emph{suggestion
meaningfulness}, and \emph{concept-axis quality}.

\paragraph{Confidence intervals.} For every metric and dimension the dashboard reports a two-sided
95\% confidence interval of the mean, computed from the per-session values under a
$t$-distribution:
%
\begin{equation}
  \text{CI}_{95} = \bar{x} \;\pm\; t_{0.975,\,n-1}\cdot \frac{s}{\sqrt{n}},
  \qquad
  s = \sqrt{\frac{1}{n-1}\sum_{i=1}^{n}\left(x_i-\bar{x}\right)^2},
  \label{eq:ci}
\end{equation}
%
where $n$ is the number of sessions contributing a value, $s$ is the Bessel-corrected sample
standard deviation, and $t_{0.975,\,n-1}$ is the critical value for $n-1$ degrees of freedom
(falling back to $1.96$ for $n\geq 31$). Bootstrapping is deliberately not used: at the sample
sizes reached here ($n\approx 30$ per condition, fewer per persona) resampling cannot recover
information the small sample lacks. Non-overlapping intervals are read as a practical signal of a
real difference between conditions, not as a formal hypothesis test.
